% รายงานโครงงานศึกษาอิสระ ITIS2026 — คอมไพล์ด้วย XeTeX (tectonic): docs/report/build.sh
\documentclass[a4paper,12pt]{report}

% ---------- หน้ากระดาษ (ตามรายงานตัวอย่าง: ขอบซ้ายกว้างสำหรับเข้าเล่ม) ----------
\usepackage[top=3.4cm,bottom=2.5cm,left=3.8cm,right=2.5cm,headheight=20pt,headsep=14pt]{geometry}

% ---------- ฟอนต์ไทย: TH Sarabun New (bundled ในเล่ม, GPL-2.0 + font exception) ----------
% มาตรฐานเอกสารราชการไทย: TH Sarabun New ที่ 16pt -- ห้ามใส่ Scale เพิ่ม ขนาดจะใหญ่เกินจริง
\usepackage{fontspec}
\setmainfont{THSarabunNew}[Path=fonts/,Extension=.ttf,
  UprightFont=*,BoldFont=*-Bold,ItalicFont=*-Italic,BoldItalicFont=*-BoldItalic,
  WordSpace={1,0,0}]
\setmonofont{DejaVu Sans Mono}[Scale=0.78]
% Thai line breaking: build/thai_zwsp.py puts U+200B between words as a break
% point only -- natural width 0, NO stretch, or justification visibly spaces
% out every Thai word (reads as garbled/broken Thai). Looser lines end with a
% little extra white space at the margin instead; that is the correct trade.
\catcode"200B=\active
\protected\def^^^^200b{\hskip0pt\relax}
% A handful of literal ASCII spaces survive inside Thai prose (clause breaks,
% a space before an inline English/CVSS token) alongside the ZWSP points
% above. Normal interword glue DOES stretch -- so on a line where ZWSP (zero
% stretch) supplies every other break, those rare real spaces end up as the
% ONLY elastic glue and absorb ALL of that line's justification stretch by
% themselves, rendering as one huge gap instead of being spread thin across
% the line (found 2026-10-21 via `pdftotext -layout` showing runs of 3+
% spaces mid-sentence). WordSpace={1,0,0} above keeps the normal space WIDTH
% but removes its stretch/shrink too, so it behaves like the ZWSP points
% around it and the same "looser lines, no visible gap" trade-off applies
% uniformly instead of concentrating on whichever real space happens to be
% on a given line.
%
% Same mechanism shows up a second way: default math spacing around a
% relation like "=" (\thickmuskip) and "," (\medmuskip) DOES stretch, so a
% Thai line that contains inline math (e.g. $\{v : \Delta(v) = \infty\}$)
% can dump its whole justification deficit into that one "=" sign instead,
% rendering as one huge gap around it. Zero that stretch too, for the same
% reason as WordSpace above.
\thickmuskip=5mu plus 0mu minus 0mu
\medmuskip=4mu plus 0mu minus 4mu

\usepackage{graphicx,xcolor,booktabs,array,multirow,longtable,float,enumitem,setspace,amsmath}
\usepackage{fancyhdr,titlesec,tocloft,caption,chngcntr}
% เลขตาราง/รูปต่อเนื่องทั้งเล่ม (รูปที่ 23, ตารางที่ 7) เหมือนรายงานตัวอย่าง
\counterwithout{table}{chapter}
\counterwithout{figure}{chapter}
\usepackage[hidelinks,bookmarks=false,pdfpagelabels=false]{hyperref}

% ---------- ขนาดตัวอักษร: เนื้อหา 16pt เหมือนรายงานตัวอย่าง ----------
% \scriptsize/\tiny were left out of this block for a while -- they kept
% the documentclass's stock 12pt-based values (8pt/6pt) instead of scaling
% up with everything else, so any \tiny/\scriptsize text (the CVSS-vector
% and edge-schema tables in the appendix) rendered at roughly half the
% size it should, and worse once \texttt's own 0.78 mono scale stacked on
% top (measured ~4.7pt/~6.2pt in the built PDF -- found 2026-10-22 from a
% user report that English words looked inconsistently sized).
\renewcommand{\tiny}{\fontsize{8pt}{10pt}\selectfont}
\renewcommand{\scriptsize}{\fontsize{10pt}{13pt}\selectfont}
\renewcommand{\footnotesize}{\fontsize{13pt}{17pt}\selectfont}
\renewcommand{\small}{\fontsize{14pt}{19pt}\selectfont}
\renewcommand{\normalsize}{\fontsize{16pt}{22pt}\selectfont}
\renewcommand{\large}{\fontsize{18pt}{24pt}\selectfont}
\renewcommand{\Large}{\fontsize{20pt}{28pt}\selectfont}
\normalsize
\setlength{\parindent}{1.25cm}
\setlength{\parskip}{2pt}
\setlength{\emergencystretch}{3em}
\sloppy
\setlist{itemsep=1pt,topsep=3pt,parsep=0pt}

% ---------- ชื่อภาษาไทย ----------
\renewcommand{\contentsname}{สารบัญ}
\renewcommand{\listfigurename}{สารบัญรูป}
\renewcommand{\listtablename}{สารบัญตาราง}
\renewcommand{\bibname}{เอกสารอ้างอิง}
\renewcommand{\figurename}{รูปที่}
\renewcommand{\tablename}{ตารางที่}
\renewcommand{\appendixname}{ภาคผนวก}

% ---------- หัวข้อ: บทอยู่กึ่งกลาง ตัวหนา; หัวข้อย่อยชิดซ้าย ----------
\titleformat{\chapter}[block]{\centering\bfseries\Large}{บทที่ \thechapter}{0.6em}{}
\titlespacing*{\chapter}{0pt}{0pt}{16pt}
\titleformat{\section}{\bfseries\large}{\thesection}{0.6em}{}
\titlespacing*{\section}{0pt}{14pt}{4pt}
\titleformat{\subsection}{\bfseries\normalsize}{\thesubsection}{0.6em}{}
\titlespacing*{\subsection}{0pt}{10pt}{3pt}
\titleformat{\subsubsection}{\bfseries\normalsize}{\thesubsubsection}{0.6em}{}
\titlespacing*{\subsubsection}{0pt}{8pt}{2pt}
% บทที่ไม่มีเลขบท (บทคัดย่อ ฯลฯ)
\newcommand{\frontchapter}[1]{\clearpage\begin{center}\bfseries\Large #1\end{center}\vspace{6pt}}

% ---------- คำบรรยาย: ตารางอยู่บนตาราง, รูปอยู่ใต้รูป ----------
\captionsetup{font=small,labelfont=bf,textfont=bf,labelsep=space,skip=6pt}
\captionsetup[table]{position=above}
\captionsetup[figure]{position=below}

% ---------- สารบัญ: บทตัวหนา มีจุดไข่ปลา ----------
\renewcommand{\cftchapfont}{\bfseries}
\renewcommand{\cftchappagefont}{\bfseries}
\renewcommand{\cftchappresnum}{บทที่ }
\setlength{\cftchapnumwidth}{4.6em}
\renewcommand{\cftchapleader}{\bfseries\cftdotfill{\cftdotsep}}
\setlength{\cftbeforechapskip}{6pt}
\setlength{\cftsecindent}{1.5em}\setlength{\cftsecnumwidth}{2.6em}
\setlength{\cftsubsecindent}{4em}\setlength{\cftsubsecnumwidth}{3.4em}
\setlength{\cftsubsubsecindent}{7.4em}\setlength{\cftsubsubsecnumwidth}{4.2em}
\setlength{\cftfigindent}{0pt}\setlength{\cftfignumwidth}{4.5em}
\setlength{\cfttabindent}{0pt}\setlength{\cfttabnumwidth}{5.2em}
\renewcommand{\cftfigpresnum}{รูปที่ }
\renewcommand{\cfttabpresnum}{ตารางที่ }
\setcounter{tocdepth}{3}
% ชื่อ สารบัญ / สารบัญตาราง / สารบัญรูป อยู่กึ่งกลาง
\renewcommand{\cfttoctitlefont}{\hfill\bfseries\Large}
\renewcommand{\cftaftertoctitle}{\hfill}
\renewcommand{\cftlottitlefont}{\hfill\bfseries\Large}
\renewcommand{\cftafterlottitle}{\hfill}
\renewcommand{\cftloftitlefont}{\hfill\bfseries\Large}
\renewcommand{\cftafterloftitle}{\hfill}

% ---------- เลขหน้า: ส่วนหน้าเป็นพยัญชนะไทย (ก ข ค …), ส่วนเนื้อหาเป็นเลขอารบิก ----------
\newcommand{\thaiconsonant}[1]{\ifcase#1\or ก\or ข\or ค\or ง\or จ\or ฉ\or ช\or ซ\or ฌ\or ญ\or ฎ\or ฏ\or ฐ\or ฑ\or ฒ\or ณ\or ด\or ต\or ถ\or ท\or ธ\or น\or บ\or ป\or ผ\or ฝ\or พ\or ฟ\or ภ\or ม\or ย\or ร\or ล\or ว\or ศ\or ษ\or ส\or ห\or ฬ\or อ\or ฮ\fi}
\pagestyle{fancy}
\fancyhf{}
\fancyhead[R]{\small\thepage}
\renewcommand{\headrulewidth}{0pt}
\fancypagestyle{plain}{\fancyhf{}\fancyhead[R]{\small\thepage}\renewcommand{\headrulewidth}{0pt}}

% ตัวช่วยสำหรับร่างที่ยังไม่เขียน
\newcommand{\draftnote}[1]{\par\noindent{\color{gray}\itshape [ร่าง — #1]}\par}

% ข้อมูลหน้าปก — แก้ที่นี่ที่เดียว
\newcommand{\TitleTH}{การหาจุดตัด (Choke Point) ของเส้นทางโจมตีโครงสร้างพื้นฐานแบบ Hybrid:\\วิธีให้น้ำหนักที่อธิบายที่มาได้ และการวัดว่าผลลัพธ์เชื่อถือได้แค่ไหน}
% TODO: ชื่ออังกฤษเป็นร่างแปลของผู้ช่วย — โปรดตรวจ/แก้
\newcommand{\TitleEN}{Choke Point Analysis of Attack Paths in Hybrid Infrastructure:\\Explainable Edge Weighting and Result Reliability}
\newcommand{\Author}{นายธนบดี อังกุลดี 67070065}
\newcommand{\AuthorShort}{ธนบดี อังกุลดี}
\newcommand{\CourseLine}{06016421 Information Technology Infrastructure Security}
\newcommand{\Advisor}{ผศ.ดร. สุเมธ ประภาวัต}
\newcommand{\Faculty}{คณะเทคโนโลยีสารสนเทศ สถาบันเทคโนโลยีพระจอมเกล้าเจ้าคุณทหารลาดกระบัง}
\newcommand{\Term}{ภาคเรียนที่ 1 ปีการศึกษา 2569}


\begin{document}
% ---------- ปก ----------
\thispagestyle{empty}
\begin{center}
\includegraphics[width=3.2cm]{figures/kmitl_logo.png}\\[0.7cm]
{\bfseries\Large รายงานโครงงานศึกษาอิสระ}\\[0.6cm]
{\bfseries \TitleEN}\\[0.5cm]
\TitleTH\\[0.9cm]
{\bfseries จัดทำโดย}\\[0.1cm]
\Author\\[0.6cm]
{\bfseries เสนอ}\\[0.1cm]
\Advisor\\[0.4cm]
\CourseLine\\[1.0cm]
{\bfseries \Faculty}\\[0.1cm]
{\bfseries \Term}
\end{center}
\clearpage
\pagenumbering{arabic}
\renewcommand{\thepage}{\thaiconsonant{\value{page}}}
\setcounter{page}{1}

% ---------- บทคัดย่อ ----------
\frontchapter{บทคัดย่อ}
เครื่องมือบริหารเส้นทางโจมตี (Attack Path Management) จัดอันดับเส้นทางที่อันตรายด้วยคะแนนที่อธิบายที่มาไม่ได้
และไม่วัดว่าอันดับเปลี่ยนเท่าไรเมื่อน้ำหนักคลาดเคลื่อน
โครงงานนี้สร้างกราฟการโจมตีของโครงสร้างพื้นฐานแบบ Hybrid ที่ทุกเส้นเชื่อม (edge) มีความน่าจะเป็น $p(e)$ อธิบายที่มาได้
(EPSS สำหรับ edge ที่มี CVE, ค่าที่ fit จากความสัมพันธ์ CVSS--EPSS สำหรับ misconfiguration
และ $p=1$ สำหรับความสัมพันธ์เชิงโครงสร้าง)
หาจุดตัดด้วย dominator tree เทียบกับ weighted node min-cut
และวัดความเสถียรของอันดับด้วย Jaccard@k และ Kendall's $\tau$

ทดลองบนกราฟจำลอง 216 โครงสร้าง (864 กราฟเมื่อรวมโมเดลน้ำหนัก 4 แบบ) และ lab ที่ใช้ Terraform, LocalStack, Cartography และ Neo4j
ผลคือ (1) อันดับทนต่อการรบกวนน้ำหนัก (Jaccard@5 เท่ากับ 0.954 ที่ $\pm 50\%$)
แต่ไม่ทนต่อ edge ที่ขาดหาย (0.49 เมื่อลบ 5\% และ 0.15 เมื่อลบ 20\%)
(2) จุดตัดที่การันตีทุกเส้นทางพบเพียง 18 จาก 216 กราฟ (8.3\%) ส่วนใหญ่ในเครือข่ายแบบ micro-segmentation
(3) weighted min-cut ที่ใช้น้ำหนัก EPSS อยู่ในเซตของ dominator tree เพียง 72\% ของกรณีที่เทียบได้
และฟังก์ชัน fit จาก CVSS อธิบายความแปรปรวนได้น้อย ($R^2$ ในสเกล logit เท่ากับ 0.219)

ข้อจำกัดคือกราฟจำลองไม่มีค่าจริงสำหรับเทียบความถูกต้อง จึงสรุปไม่ได้ว่า EPSS แม่นกว่า CVSS
และ lab ใช้ Docker, nmap และ LocalStack แทน GNS3, Nessus และ AWS จริง

\vspace{0.4cm}
\noindent\textbf{คำสำคัญ:} จุดตัด (Choke Point), กราฟการโจมตี, Dominator Tree, EPSS, CVSS, โครงสร้างพื้นฐานแบบ Hybrid


% ---------- กิตติกรรมประกาศ ----------
\frontchapter{กิตติกรรมประกาศ}
\setlength{\parindent}{1.25cm}
โครงงานนี้สำเร็จลุล่วงได้ด้วยความกรุณาจาก \Advisor{} อาจารย์ที่ปรึกษา ที่ให้คำแนะนำและข้อเสนอแนะตลอดการทำโครงงาน
ผู้จัดทำขอขอบพระคุณเป็นอย่างสูงมา ณ โอกาสนี้

ขอขอบคุณคณะเทคโนโลยีสารสนเทศ สถาบันเทคโนโลยีพระจอมเกล้าเจ้าคุณทหารลาดกระบัง
ที่ให้ความรู้และสภาพแวดล้อมในการศึกษา รวมถึงครอบครัวและเพื่อน ๆ ที่ให้กำลังใจตลอดการทำโครงงาน

\vspace{1.2cm}
\hfill \AuthorShort

% ---------- สารบัญ ----------
\clearpage
\tableofcontents
\clearpage
\listoftables
\clearpage
\listoffigures


% ================= เนื้อหา =================
\clearpage\pagenumbering{arabic}
\chapter{บทนำ}
\section{ที่มาและความสำคัญของปัญหา}
\draftnote{องค์กรไฮบริด, Attack Path Management (BloodHound/Wiz/XM Cyber), edge ข้ามโดเมน เช่น \texttt{\~{}/.aws/credentials} และ SSRF ไป IMDS}
\section{ปัญหาที่ต้องการแก้ไข}
\draftnote{น้ำหนัก edge ขาดความน่าเชื่อถือ ยังไม่มีใครวัดว่าอันดับเปลี่ยนเท่าไรเมื่อน้ำหนักผิด}
\section{วัตถุประสงค์}
\draftnote{5 ข้อตามข้อเสนอโครงงาน}
\section{ขอบเขตของโครงงาน}
\draftnote{ในขอบเขต/นอกขอบเขต และสมมติฐาน OR-semantics, monotonic, edge independence}
\section{ประโยชน์ที่คาดว่าจะได้รับ}
\draftnote{5 ข้อตามข้อเสนอโครงงาน}

\chapter{ทฤษฎีและงานที่เกี่ยวข้อง}
\section{กราฟการโจมตี (Attack Graph)}\draftnote{Sheyner 2002; MulVAL 2005}
\section{การเสริมความมั่นคงต้นทุนต่ำสุด}\draftnote{Noel 2003}
\section{Dominator Tree}\draftnote{Lengauer--Tarjan 1979}
\section{CVSS และ EPSS}\draftnote{FIRST; EPSS วัดโอกาสถูกโจมตีที่ใดที่หนึ่ง ไม่ใช่โอกาสสำเร็จ}
\section{เครื่องมือที่เกี่ยวข้อง}\draftnote{BloodHound, Wiz, XM Cyber, Cartography, Neo4j}
\section{ช่องว่างที่โครงงานนี้เติมเต็ม}\draftnote{}

\chapter{การพัฒนาระบบ}
\section{โครงสร้างโปรเจกต์}\draftnote{}
\section{ท่อประมวลผลข้อมูลและการ fit}\draftnote{EPSS + NVD 46,360 CVE}
\section{ตัววิเคราะห์จุดตัด}\draftnote{}
\section{การทำซ้ำผลได้ (Reproducibility)}\draftnote{seed, ไม่ใช้ \texttt{hash()}, sorted set, Writer ไม่ต่อท้ายไฟล์}
\section{Lab: Terraform, Cartography, Neo4j}\draftnote{กฎสองแหล่งอิสระต้องตรงกัน และ tamper test}
\section{การทดสอบซอฟต์แวร์}\draftnote{70 tests}

\chapter{การทดลอง ประเมินผล และอภิปรายผล}

\section{ผลการประมาณค่า $p(e)$ จาก CVSS (วัตถุประสงค์ข้อ 2)}

edge ที่เป็น misconfiguration ไม่มี CVE จึงไม่มีคะแนน EPSS ให้ใช้โดยตรง
โครงงานนี้จึงกำหนดเวกเตอร์ CVSS ให้แต่ละคลาสของ misconfiguration ตาม rubric (ภาคผนวก ก)
แล้วใช้ฟังก์ชันที่เรียนรู้ความสัมพันธ์ระหว่างเวกเตอร์ CVSS กับคะแนน EPSS ของ CVE จริง
แปลงเวกเตอร์นั้นเป็นค่าความน่าจะเป็น $p(e)$
หัวข้อนี้รายงานว่าฟังก์ชันดังกล่าวแม่นเพียงใด
เพราะความคลาดเคลื่อนของ $p(e)$ ส่งผลต่ออันดับเส้นทางโดยตรง

\subsection{ข้อมูลและวิธีการ}

ข้อมูลที่ใช้เป็น CVE จำนวน 46{,}360 รายการที่มีทั้งเวกเตอร์ CVSS (จาก NVD จำนวน 33 หน้า)
และคะแนน EPSS จาก snapshot วันที่ 2026-08-26 โดยเป็น CVE ที่เผยแพร่ระหว่าง 2017-03-20 ถึง 2026-08-06
การสุ่มทุกขั้นตอนใช้ seed 20260826
ตัวแปรตามคือ EPSS ในสเกล logit
และเปรียบเทียบ 4 โมเดล ได้แก่ ค่าเฉลี่ยคงที่ (constant baseline),
การถดถอยเชิงเส้นจากคะแนนฐาน CVSS เพียงตัวเดียว,
การถดถอยเชิงเส้นจาก 8 metric ของ CVSS (OLS)
และ gradient boosting จาก 8 metric เดียวกัน (GBM)
ตัวชี้วัดความแม่นในตารางที่~\ref{tab:fit} คำนวณบนข้อมูลทดสอบ 25\%

\subsection{ผลการ fit}

\begin{table}[H]
\centering
\caption{ผลการ fit จากเวกเตอร์ CVSS ไปเป็น EPSS (ข้อมูลทดสอบ 25\%)}\label{tab:fit}
\footnotesize\setlength{\tabcolsep}{4pt}
\begin{tabular}{p{5.6cm}rrrr}
\toprule
โมเดล & $R^2$ (logit) & $R^2$ (prob) & $R^2$ (top 10\%) & $\rho$ \\
\midrule
constant baseline      & 0.000 & $-0.032$ & $-0.452$ & --- \\
OLS จากคะแนนฐานตัวเดียว & 0.118 & $-0.018$ & $-0.432$ & 0.354 \\
\textbf{OLS จาก 8 metric (ใช้จริง)} & \textbf{0.219} & $-0.010$ & $-0.420$ & 0.548 \\
GBM จาก 8 metric (ไม่ได้ใช้) & 0.268 & $-0.005$ & $-0.415$ & 0.601 \\
\bottomrule
\end{tabular}
\end{table}

ผลลัพธ์สำคัญมี 3 ข้อ
ข้อแรก โมเดลที่ใช้ตั้งราคา edge คือ OLS จาก 8 metric
เพราะสัมประสิทธิ์แต่ละตัวบวกกันในสเกล logit จึงตรวจสอบย้อนกลับได้ทีละ metric
ซึ่งตรงกับวัตถุประสงค์ว่า $p(e)$ ต้องอธิบายที่มาได้
GBM ให้ค่าสูงกว่าเล็กน้อยแต่เป็นกล่องดำและไม่ได้นำไปใช้ ตัวเลขทั้งหมดในรายงานจึงเป็นของ OLS
ข้อสอง $R^2$ ในสเกลความน่าจะเป็นติดลบในทุกโมเดล
หมายความว่าในหน่วยที่ $p(e)$ ถูกใช้จริง ฟังก์ชันที่ fit ได้ไม่ดีกว่าการทายค่าเฉลี่ย
และอธิบายความแปรปรวนได้เพียงประมาณหนึ่งในห้า (0.219) เฉพาะในสเกล logit
ข้อสาม ความคลาดเคลื่อนไม่ได้เป็นสัญญาณรบกวนเท่ากันทุกช่วง
แต่เกิดจากการที่ฟังก์ชัน ``ไปไม่ถึงหาง'' ของการแจกแจง
ค่าที่ทำนายสูงสุดคือ 0.0186 ขณะที่ EPSS จริงสูงสุดถึง 0.99999
ฟังก์ชันจึงถดถอยเข้าหาค่าฐาน (base rate) ประมาณ 0.0055

ค่า $R^2 = 0.219$ ในตารางที่~\ref{tab:fit} ใช้การแบ่งข้อมูลทดสอบแบบสุ่มราย CVE
ซึ่งมองโลกในแง่ดีกว่าการใช้งานจริงเล็กน้อย เพราะ CVE ที่คล้ายกันอาจตกอยู่คนละฝั่งของการแบ่ง
เมื่อกันทั้งหน้า NVD ออกจากกัน $R^2$ ลดเหลือ 0.156 และเมื่อกันทั้งปีออกเหลือ 0.150
กรณีที่เข้มที่สุดคือฝึกด้วย CVE ถึงปี 2022 แล้วทดสอบกับปี 2023 เป็นต้นไป $R^2$ ติดลบ ($-0.863$)
แต่ Spearman $\rho$ ยังคงอยู่ที่ 0.542 แปลว่า \emph{ระดับ} ของ $p(e)$ เลื่อนไปตามเวลา
ขณะที่ \emph{อันดับ} ยังรักษาไว้ได้ (ตรวจด้วย \texttt{scripts/audit\_data.py})
จึงไม่ควรอ้างว่าฟังก์ชันนี้ทำนายค่า EPSS ของ CVE ใหม่ได้ ใช้ได้เพียงเพื่อจัดอันดับ

\subsection{ความคลาดเคลื่อนในหน่วย log}

เนื่องจากต้นทุนของเส้นทางคือผลรวมของ $-\log p$ บน edge
ค่าที่กำหนดว่าอันดับเส้นทางเปลี่ยนหรือไม่คือการกระจายของ
$r = \log p_{\text{จริง}} - \log p_{\text{ทำนาย}}$
ไม่ใช่ค่าสหสัมพันธ์ $\rho$ เพราะการแปลงแบบ monotone ต่อ edge สามารถสลับอันดับของเส้นทางได้
(เส้นทางที่มี 2 edge ที่ $p=0.9$ ได้ $0.81$ ชนะเส้นทางที่มี 1 edge ที่ $p=0.8$
แต่ถ้า map $0.9 \mapsto 0.5$ และ $0.8 \mapsto 0.45$ อันดับจะกลับกัน)
อันดับของเส้นทางคงเดิมพอดีเมื่อฟังก์ชันเป็น power law $\hat{p} \propto p^{k}$
จึงนำค่ามัธยฐานของ $r$ ออก (ซึ่งคือค่า $k$) แล้ววัดความกว้างที่เหลือ
ตารางที่~\ref{tab:resid} รายงานค่าสัมบูรณ์ $|r|$ จากการทำนายแบบ 5-fold out-of-fold บน CVE ทั้ง 46{,}360 รายการ
โดยแยกกลุ่มที่ EPSS สูงที่สุด 10\% ออกมาต่างหาก เพราะ edge ที่ตัดสินเส้นทางคือ edge ที่มีโอกาสสูง

\begin{table}[H]
\centering
\caption{ความคลาดเคลื่อน $|r|$ ในหน่วย nats แยกตามช่วงของ EPSS จริง}\label{tab:resid}
\small
\begin{tabular}{lrrr}
\toprule
กลุ่ม & $n$ & $|r|$ มัธยฐาน & $|r|$ เปอร์เซ็นไทล์ที่ 90 \\
\midrule
ทั้งหมด (pooled) & 46{,}360 & 0.587 & 1.483 \\
ล่าง 90\% & 41{,}724 & 0.523 & 1.172 \\
\textbf{EPSS สูงสุด 10\%} & 4{,}636 & \textbf{1.938} & \textbf{3.980} \\
\bottomrule
\end{tabular}
\end{table}

กลุ่มที่ EPSS สูงสุด 10\% คลาดเคลื่อนมากกว่ากลุ่มล่าง 90\% ประมาณ 3.7 เท่า (1.938 เทียบกับ 0.523)
และการรายงานค่า pooled เพียงค่าเดียว (0.587) จะทำให้ฟังก์ชันดูแม่นเกินจริง
ค่า 1.938 nats เทียบเท่าตัวคูณ $e^{1.938} \approx 6.95$ เท่าบน $p(e)$
ขณะที่แถบการรบกวน $\pm 10$--$50\%$ ในข้อเสนอโครงงานมีขนาดมากสุดเพียง $\log 1.5 \approx 0.405$ nats (1.5 เท่า)
แถบที่กำหนดไว้จึงแคบกว่าที่วัดได้ประมาณ 4.8 เท่าที่มัธยฐาน และ 9.8 เท่าที่เปอร์เซ็นไทล์ที่ 90
(ในหน่วย log) ข้อสรุปที่ถูกต้องจึงเป็น ``แคบกว่า 4--10 เท่า แล้วแต่เปอร์เซ็นไทล์''
ไม่ใช่ ``หนึ่ง order of magnitude'' ซึ่งเป็นจริงเฉพาะที่เปอร์เซ็นไทล์ที่ 90

\subsection{ผลต่อความสามารถในการแยกอันดับเส้นทาง}

เมื่อ $|r|\le\delta$ ทุก edge เส้นทางสองเส้นที่มีความยาว $|P|$ และ $|Q|$ จะรักษาอันดับจริงไว้ได้
ก็ต่อเมื่อผลต่างต้นทุนจริงมากกว่า $\delta\,(|P|+|Q|)$
ตารางที่~\ref{tab:order} แสดงขั้นต่ำของผลต่างนั้น เมื่อใช้ $\delta$ เป็นมัธยฐานและเปอร์เซ็นไทล์ที่ 90

\begin{table}[H]
\centering
\caption{ผลต่างต้นทุนขั้นต่ำ (nats) ที่ทำให้สองเส้นทางยาวเท่ากันยังรักษาอันดับไว้ได้}\label{tab:order}
\small
\begin{tabular}{rrrrr}
\toprule
& \multicolumn{2}{c}{ค่า pooled} & \multicolumn{2}{c}{EPSS สูงสุด 10\%} \\
\cmidrule(lr){2-3}\cmidrule(lr){4-5}
ความยาวเส้นทาง & มัธยฐาน & P90 & มัธยฐาน & P90 \\
\midrule
2 edge & 2.35 & 5.93 & 7.75 & 15.92 \\
4 edge & 4.70 & 11.86 & 15.51 & 31.84 \\
8 edge & 9.39 & 23.72 & 31.01 & 63.68 \\
\bottomrule
\end{tabular}
\end{table}

สำหรับเส้นทางยาว 4 edge ผลต่างต้นทุน 4.70 nats เทียบเท่าความน่าจะเป็นจริงที่ต่างกัน
$e^{4.70} \approx 109$ เท่า (edge คุณภาพมัธยฐาน)
และ 15.51 nats ในกลุ่มหาง เทียบเท่าประมาณ 5.4 ล้านเท่า
ดังนั้นน้ำหนักจาก fit แยกอันดับเส้นทางได้เฉพาะคู่ที่ต่างกันมากเท่านั้น
และแทบแยกไม่ได้เลยเมื่อ edge ที่ตัดสินอยู่ในกลุ่มที่ EPSS สูง
เงื่อนไขนี้เป็นเงื่อนไขจำเป็น ไม่ใช่เงื่อนไขเพียงพอ
ผลที่เกิดจริงต่ออันดับจุดตัดถูกวัดโดยตรงในหัวข้อ~\ref{sec:exp6}

\section{การเปรียบเทียบ Dominator Tree กับ Weighted Min-Cut (วัตถุประสงค์ข้อ 3)}

วัตถุประสงค์ข้อ 3 ถามว่าจุดตัดสองนิยามให้คำตอบตรงกันหรือไม่ นิยามแรกคือ post-dominator
ซึ่งเป็นโครงสร้างล้วน ไม่อ่านน้ำหนัก และคืนเซตของ node ที่อยู่บน\emph{ทุก}เส้นทาง
นิยามที่สองคือ weighted node min-cut ซึ่งคืนชุด node ที่ตัดแล้วผู้โจมตีไปไม่ถึงเป้าหมาย
โดยมีต้นทุนรวมต่ำสุด การเทียบสองนิยามนี้ทำบนกราฟชุดเดียวกันทั้ง 864 กราฟถ่วงน้ำหนัก

\subsection{สามนิยามของจุดตัด}

ก่อนเทียบกับ min-cut ต้องแยกให้ชัดว่า ``จุดตัด'' หมายถึงนิยามใด เพราะทั้งสามให้เซตที่ต่างกัน

\begin{table}[H]
\centering
\caption{สามนิยามของจุดตัดบน 864 กราฟถ่วงน้ำหนัก}\label{tab:defs}
\small
\begin{tabular}{p{4.6cm}rr}
\toprule
นิยาม & กราฟที่เซตไม่ว่าง & ขนาดเซตเฉลี่ย \\
\midrule
global (ทุก entry ถึงทุก target) & 8.3\% & 4.8 \\
per-target (ทุก entry ถึง target เดียว) & 100\% & 5.8 \\
per-pair (หนึ่งคู่ entry--target) & 100\% & 9.2 \\
\bottomrule
\end{tabular}
\end{table}

นิยาม global เข้มที่สุด จุดตัดต้องอยู่บนทุกเส้นทางของทุกคู่ จึงพบเพียง 18 จาก 216 โครงสร้าง (8.3\%)
ส่วน per-target และ per-pair ไม่ว่าง 100\% ซึ่ง \textbf{เป็น artifact ของตัวสร้างกราฟ ไม่ใช่ผลต่อเครือข่ายจริง}
ในกราฟที่สร้างขึ้น target แต่ละตัวมี in-edge เพียงเส้นเดียวโดยโครงสร้าง node ต้นทางของ edge นั้นจึง
ครอบทุกเส้นทางที่เข้าถึง target เสมอ ค่า 100\% จึงสะท้อนวิธีสร้างกราฟ ไม่ใช่คุณสมบัติของเครือข่ายจริง
รายงานนี้จึงใช้นิยาม global เป็นพาดหัว และกำกับข้อจำกัดนี้ไว้ทุกที่ที่อ้างถึง per-target/per-pair

\subsection{จุดตัดเชิงโครงสร้างเป็นกรณีสุดขั้วของชั้นความน่าจะเป็น}

post-dominator set พิสูจน์ได้ว่าเท่ากับ $\{v : \Delta(v) = \infty\}$ นั่นคือ node ที่ลบแล้ว
ผู้โจมตีไปไม่ถึงเป้าหมายเลย (ต้นทุนเพิ่มเป็นอนันต์) จุดตัดเชิงโครงสร้างจึงไม่ใช่วิธีที่แยกจาก
ชั้นความน่าจะเป็น แต่เป็นกรณีสุดขั้วของมัน คุณสมบัตินี้ตรวจด้วย brute force บนกราฟสุ่ม 300 กราฟ
(\texttt{tests/test\_bruteforce.py}): จำนวน node ที่ $\Delta(v)=\infty$ เท่ากับขนาดของ global set ทุกกราฟ

\subsection{ความตรงกันของ dominator กับ min-cut}

การเทียบทำบนกราฟ 18 โครงสร้างที่ global set ไม่ว่าง (เซตว่างเทียบไม่ได้ และ \texttt{disagreement}
บนเซตว่างมีค่า 1 โดยปริยาย ไม่ได้บอกอะไรเกี่ยวกับสองนิยาม) min-cut รายงานสาม capacity ตามการตัดสินใจ
ตารางที่~\ref{tab:cut} แสดงสัดส่วนที่ node ใน min-cut เป็นสมาชิกของ dominator set ด้วย

\begin{table}[H]
\centering
\caption{สัดส่วน node ใน min-cut ที่อยู่ใน dominator set (18 โครงสร้างที่ global ไม่ว่าง)}\label{tab:cut}
\small
\begin{tabular}{lrp{3.2cm}l}
\toprule
capacity & ขนาด cut เฉลี่ย & cut อยู่ใน dominator set & บทบาท \\
\midrule
unit & 1.0 & 100\% & หน่วยตรงกับ dominator \\
prob & 1.1 & 72\% EPSS, 100\% อื่น & ค่าตั้งต้น \\
cost & 3.0 & 10\% & degenerate \\
\bottomrule
\end{tabular}
\end{table}

ผลสำคัญมีสามข้อ ข้อแรก capacity แบบ \texttt{unit} ซึ่งนับจำนวน node ที่ต้องตัด ตรงกับ dominator set
100\% เป็นผลที่ควรเป็น เพราะทั้งคู่วัดในหน่วยเดียวกันและ dominator ไม่อ่านน้ำหนัก
ข้อสอง capacity แบบ \texttt{prob} ซึ่งเป็นค่าตั้งต้น ตรงกับ dominator 100\% เมื่อใช้น้ำหนัก CVSS หรือ uniform
แต่ลดเหลือ \textbf{72\% เมื่อใช้น้ำหนัก EPSS} กล่าวคือใน 28\% ของกรณี min-cut ที่ถ่วงด้วย EPSS
เลือก node ที่ไม่ใช่ dominator เพราะ node นั้นเป็นคอขวดของ\emph{ความน่าจะเป็น} แม้ไม่ได้อยู่บนทุกเส้นทาง
นี่คือข้อค้นพบหลักของวัตถุประสงค์ข้อ 3 ที่น้ำหนักเชิงประจักษ์ให้คำตอบต่างจากโครงสร้างล้วน
ข้อสาม capacity แบบ \texttt{cost} ($-\log p$) เสื่อม (degenerate) เพราะ edge เชิงโครงสร้างมี $p=1$ จึง $w=0$
node ที่เข้าถึงได้ฟรีจึงถูกตัดด้วยต้นทุนศูนย์ min-cut ยุบไปที่เซตต้นทุนศูนย์ที่เลือกโดยพลการ
ยืนยันด้วยสัดส่วน cut ที่มีต้นทุนรวมเป็นศูนย์พอดีถึง 77\% ของแถว จึงรายงาน \texttt{cost} ไว้เป็น
คอลัมน์ sensitivity เท่านั้น ไม่ใช่ค่าหลัก

\subsection{ระดับการแบ่งส่วนเครือข่ายกับการเกิดจุดตัด}

ข้อเสนอโครงงานถามว่าเครือข่ายต้องแบ่งส่วนละเอียดแค่ไหนจุดตัดที่การันตีได้จึงจะปรากฏ
คำถามนี้ตอบได้จากสัดส่วนที่ global set ไม่ว่าง แยกตามระดับการแบ่งส่วน รายงานเป็นจำนวนกราฟ
พร้อมช่วงความเชื่อมั่น 95\% แบบ Wilson (ไม่ใช่เปอร์เซ็นต์ลอย ๆ เพราะจำนวนน้อย)

\begin{table}[H]
\centering
\caption{สัดส่วนกราฟที่มี global choke point แยกตามการแบ่งส่วน (72 กราฟต่อระดับ)}\label{tab:seg}
\small
\begin{tabular}{lrl}
\toprule
การแบ่งส่วน & กราฟที่มีจุดตัด & ช่วงความเชื่อมั่น 95\% (Wilson) \\
\midrule
flat (ไม่แบ่ง) & 1 / 72 & 0.2--7.5\% \\
3-tier & 1 / 72 & 0.2--7.5\% \\
micro-segmentation & 16 / 72 & 14--33\% \\
\midrule
รวม & 18 / 216 & 5.3--12.8\% \\
\bottomrule
\end{tabular}
\end{table}

micro-segmentation ให้จุดตัดที่การันตีได้บ่อยกว่าอีกสองแบบอย่างมีนัยสำคัญ (ช่วงความเชื่อมั่นไม่ซ้อนทับ)
ซึ่งสอดคล้องกับสัญชาตญาณ เครือข่ายที่แบ่งส่วนละเอียดบีบให้เส้นทางลอดผ่าน node เชื่อมไม่กี่ตัว
node เหล่านั้นจึงกลายเป็นคอขวดที่ทุกเส้นทางต้องผ่าน ส่วน flat และ 3-tier มีเส้นทางสำรองมากเกินกว่าจะมี
จุดตัดร่วม อย่างไรก็ตามแม้แต่ micro ก็ยังมีจุดตัดเพียง 22\% ของกราฟ แปลว่า\textbf{การพึ่งจุดตัดที่การันตีได้
เป็นกลยุทธ์ที่ใช้ได้เฉพาะส่วนน้อยของกราฟ} ส่วนที่เหลือต้องใช้การจัดอันดับ $\Delta(v)$ แทนการหาเซตที่การันตี

หมายเหตุด้านความแปรปรวน: grid รุ่นก่อน (135 โครงสร้าง, seed คนละชุด) เคยให้ค่า micro ที่ 44\%
ค่าที่ต่างกันมากนี้ย้ำว่าสัดส่วนดังกล่าวมีความแปรปรวนสูงเมื่อจำนวนกราฟที่เข้าเงื่อนไขมีน้อย
จึงต้องรายงานเป็นจำนวนกราฟพร้อมช่วงความเชื่อมั่นเสมอ ไม่ใช่ตัวเลขจุดเดียว
\section{ความเสถียรของอันดับเมื่อน้ำหนักถูกรบกวน (วัตถุประสงค์ข้อ 4)}

วัตถุประสงค์ข้อ 4 ถามว่าอันดับจุดตัดขยับแค่ไหนเมื่อน้ำหนัก edge คลาดเคลื่อน การทดลองนี้ใช้สอง
รูปแบบของสัญญาณรบกวน (noise) คือแบบสมมาตร $\pm 10$--$50\%$ ตามที่ข้อเสนอโครงงานระบุ
และแบบเชิงประจักษ์ (empirical) ซึ่งสุ่มจาก residual จริงของการ fit ในหัวข้อ~\ref{sec:exp6}
คูณด้วยสเกล 0.5 ถึง 8 เท่า (สเกล $\geq 4$ เพิ่มเข้ามาเพื่อหาจุดที่อันดับเริ่มพัง เกินกว่าที่ข้อเสนอกำหนด)
วัดด้วย Jaccard@5 เทียบอันดับก่อนและหลังรบกวน บนกราฟเดิมและโมเดลน้ำหนักเดิม

\begin{table}[H]
\centering
\caption{Jaccard@5 เฉลี่ยเมื่อรบกวนน้ำหนัก (ทุกโมเดลน้ำหนักรวมกัน)}\label{tab:noise}
\small
\begin{tabular}{lrrrrr}
\toprule
& \multicolumn{4}{c}{แบบสมมาตร ($\pm$)} & \\
\cmidrule(lr){2-5}
& 10\% & 20\% & 30\% & 50\% & \\
Jaccard@5 & 0.979 & 0.968 & 0.964 & 0.954 & \\
\midrule
& \multicolumn{4}{c}{แบบเชิงประจักษ์ (เท่าของ residual จริง)} & \\
\cmidrule(lr){2-5}
& 0.5$\times$ & 1$\times$ & 2$\times$ & 4$\times$ & 8$\times$ \\
Jaccard@5 & 0.968 & 0.952 & 0.935 & 0.920 & 0.912 \\
\bottomrule
\end{tabular}
\end{table}

ผลที่เด่นชัดคือ \textbf{ไม่พบจุดที่อันดับพังแม้ขยายสัญญาณรบกวนไปถึง 8 เท่าของ residual จริง}
Jaccard@5 ลดลงอย่างช้า ๆ จาก 0.979 ที่ $\pm10\%$ เหลือ 0.954 ที่ $\pm50\%$ และเหลือ 0.912
ที่ความคลาดเคลื่อนใหญ่กว่าที่วัดได้จริงถึง 8 เท่า ซึ่งสอดคล้องกับผลในหัวข้อ~\ref{sec:exp6}
ว่าอันดับ node ทนทานกว่าที่ fit error เพียงอย่างเดียวจะทำนายได้

ตารางที่~\ref{tab:noise_model} แยกตามโมเดลน้ำหนัก เพื่อตรวจว่าความทนทานเป็นคุณสมบัติของ
โมเดลใดโมเดลหนึ่งหรือของโครงสร้างกราฟเอง

\begin{table}[H]
\centering
\caption{Jaccard@5 แยกตามโมเดลน้ำหนัก (แบบสมมาตร $\pm50\%$ และเชิงประจักษ์ $8\times$)}\label{tab:noise_model}
\small
\begin{tabular}{lrrrr}
\toprule
& epss & cvss & uniform & random (null) \\
\midrule
สมมาตร $\pm50\%$ & 0.986 & 0.909 & 0.966 & 0.956 \\
เชิงประจักษ์ $8\times$ & 0.921 & 0.908 & 0.914 & 0.905 \\
\bottomrule
\end{tabular}
\end{table}

ที่สำคัญคือ \textbf{แม้แต่โมเดล \texttt{random} ซึ่งสุ่มน้ำหนักทิ้งทั้งหมด} ก็ยังให้ Jaccard@5
สูงถึง 0.905--0.956 เมื่อถูกรบกวนเพิ่ม นั่นหมายความว่าความทนทานที่สังเกตได้ไม่ได้มาจาก
คุณภาพของน้ำหนักที่ใช้ แต่มาจาก\textbf{โครงสร้างของกราฟเอง} — มีจำนวน node จำกัดที่อยู่บน
เส้นทางถูกสุดได้ และการรบกวนน้ำหนักเพียงเล็กน้อยไม่พอจะเปลี่ยนว่า node ใดอยู่ในกลุ่มนั้น
ข้อสรุปนี้ต้องอ่านคู่กับหัวขัดถัดไป ซึ่งแสดงว่าอันดับไม่ทนทานต่อ\emph{โครงสร้าง}ที่เปลี่ยน

\section{ความเสถียรของอันดับเมื่อ edge ขาดหาย}

การทดลองนี้ลบ edge แบบสุ่ม 5, 10 และ 20\% ของกราฟ เพื่อจำลองกรณีที่ข้อมูลเก็บมาไม่ครบ
ต่างจากการรบกวนน้ำหนักในหัวข้อก่อน การลบ edge เปลี่ยน\emph{โครงสร้าง}ของกราฟโดยตรง
และอาจตัดบางคู่ (entry, target) ออกจากกันทั้งคู่

\begin{table}[H]
\centering
\caption{ผลของการลบ edge แบบสุ่ม}\label{tab:drop}
\small
\begin{tabular}{rrrrp{2.6cm}}
\toprule
edge ที่ลบ & Jaccard@5 ดิบ & Kendall $\tau$ & คู่ที่ยังถึงกัน & Jaccard@5 เงื่อนไข (n กราฟ) \\
\midrule
5\%  & 0.493 & 0.672 & 69\% & 0.749 (177) \\
10\% & 0.321 & 0.526 & 47\% & 0.650 (97) \\
20\% & 0.145 & 0.371 & 20\% & 0.460 (17) \\
\bottomrule
\end{tabular}
\end{table}

คอลัมน์ ``ดิบ'' คือ Jaccard@5 เฉลี่ยทุกกราฟ ซึ่งปนผลสองอย่างเข้าด้วยกัน: อันดับเปลี่ยนเพราะ
ต้นทุนเส้นทางเปลี่ยน และคู่ (entry, target) บางคู่หลุดจากกันไปเลยทำให้ node ที่เคยอยู่ในอันดับ
หายไปโดยสิ้นเชิง เพื่อแยกสองผลนี้ คอลัมน์ ``เงื่อนไข'' คำนวณเฉพาะกราฟที่ยังคงคู่ (entry, target)
ไว้ได้อย่างน้อย 90\% หลังลบ edge แม้เงื่อนไขนี้จะทำให้ Jaccard@5 สูงขึ้นมาก (0.749 เทียบ 0.493
ที่ 5\%) แต่ก็ยังลดลงต่อเนื่องเมื่อลบมากขึ้น และจำนวนกราฟที่ผ่านเงื่อนไขลดฮวบเหลือเพียง 17 กราฟ
ที่ 20\% ทำให้ตัวเลขนั้นมีความแปรปรวนสูง (ช่วงความเชื่อมั่น 95\% กว้างถึง 0.32--0.66)

ผลสำคัญคือ \textbf{edge ที่ขาดหายกระทบอันดับรุนแรงกว่าน้ำหนักที่คลาดเคลื่อนมาก} ลบเพียง 5\%
ของ edge ลด Jaccard@5 ลงมากกว่าการรบกวนน้ำหนักถึง 8 เท่าของ residual จริงเสียอีก
(0.493 เทียบ 0.912) ข้อสรุปเชิงปฏิบัติคือ ความสมบูรณ์ของข้อมูล (การเก็บ edge ให้ครบ)
สำคัญกว่าความแม่นยำของน้ำหนักบน edge ที่เก็บมาได้แล้ว

\section{การเปรียบเทียบโมเดลน้ำหนัก 4 แบบ}

หัวข้อนี้เทียบอันดับจุดตัดที่ได้จากน้ำหนัก 4 แบบบนกราฟโครงสร้างเดียวกัน (paired comparison):
EPSS, CVSS (คะแนนฐานหารสิบ), uniform ($p=0.5$ ทุก edge ที่ไม่ใช่ deterministic) และ random
(สุ่ม $p$ ใหม่ทุกครั้ง) โมเดลหลังสุดเป็น null model: โครงสร้างใดก็ตามที่รอดจากการสุ่มทิ้งนี้
คือโครงสร้างจริง ไม่ใช่สัญญาณจากน้ำหนัก

\begin{table}[H]
\centering
\caption{Jaccard@5 ระหว่างคู่โมเดลน้ำหนัก (216 กราฟ จับคู่กัน)}\label{tab:models}
\small
\begin{tabular}{lrr}
\toprule
คู่โมเดล & Jaccard@5 & Kendall $\tau$ \\
\midrule
epss -- uniform & 0.974 & 0.947 \\
epss -- cvss & 0.951 & 0.884 \\
epss -- random (null) & 0.924 & 0.834 \\
uniform -- random (null) & 0.929 & 0.832 \\
cvss -- uniform & 0.941 & 0.878 \\
cvss -- random (null) & 0.910 & 0.817 \\
\bottomrule
\end{tabular}
\end{table}

ตัวเลขที่ต้องอ่านคู่กันคือแถวที่สองและสาม: EPSS เทียบกับ CVSS ให้ Jaccard@5 เท่ากับ 0.951
ขณะที่ EPSS เทียบกับ \emph{random} ให้ 0.924 — ต่างกันเพียง 0.027 \textbf{จึงไม่มีหลักฐานเพียงพอ
ที่จะอ้างว่า EPSS ให้อันดับที่ ``ถูกต้องกว่า'' CVSS} เพราะกราฟจำลองไม่มีค่าจริง (ground truth)
ให้เทียบ ตัวเลขทั้งหมดในตารางนี้วัดได้แค่ \emph{ความต่างกันของอันดับ} ระหว่างสองน้ำหนัก
ไม่ใช่ความถูกต้อง การที่ EPSS ใกล้ CVSS มากกว่าใกล้ random เพียงเล็กน้อย บ่งชี้ว่าทั้งสองโมเดล
ที่มีข้อมูลจริงรองรับให้อันดับที่คล้ายกันในระดับหนึ่ง แต่ความคล้ายนั้นเทียบกับ null model แล้ว
มีขนาดเล็ก สรุปว่าการทดลองนี้วัดได้เพียง disagreement rate และ
ความเสถียรของอันดับ ไม่ใช่ความแม่นยำ
\section{ผลของความคลาดเคลื่อนจากการ fit ต่ออันดับจุดตัด}\label{sec:exp6}\draftnote{§5.7 — exp 6}
\section{สัดส่วนเส้นทางที่ข้ามโดเมน}\draftnote{§5.8 — exp 7}
\section{ผลจาก lab จริง (วัตถุประสงค์ข้อ 5)}\draftnote{§5.9}
\section{อภิปรายผลและภัยต่อความตรงของผล}\draftnote{§5.10}
\section{ข้อจำกัด}\draftnote{§5.11}

\chapter{สรุปผลและแนวทางพัฒนา}
\section{สรุปผลตามวัตถุประสงค์}\draftnote{}
\section{ข้อจำกัด}\draftnote{}
\section{แนวทางพัฒนาต่อ}\draftnote{}


\begin{thebibliography}{99}
\addcontentsline{toc}{chapter}{เอกสารอ้างอิง}
\bibitem{sheyner2002} O. Sheyner et al., ``Automated Generation and Analysis of Attack Graphs,'' \emph{IEEE S\&P}, 2002.
\bibitem{ou2005} X. Ou et al., ``MulVAL: A Logic-based Network Security Analyzer,'' \emph{USENIX Security}, 2005.
\bibitem{noel2003} S. Noel et al., ``Efficient Minimum-Cost Network Hardening via Exploit Dependency Graphs,'' \emph{ACSAC}, 2003.
\bibitem{lengauer1979} T. Lengauer and R. Tarjan, ``A Fast Algorithm for Finding Dominators in a Flowgraph,'' \emph{ACM TOPLAS}, 1979.
\bibitem{jacobs2019} J. Jacobs, S. Romanosky, B. Edwards, M. Roytman, and I. Adjerid, ``Exploit Prediction Scoring System (EPSS),'' presented at Black Hat USA, 2019. [Online]. Available: \url{https://arxiv.org/abs/1908.04856}
\bibitem{cvss31} FIRST, ``Common Vulnerability Scoring System version 3.1: Specification Document.'' [Online]. Available: \url{https://www.first.org/cvss/v3.1/specification-document}
\bibitem{epss} FIRST, ``Exploit Prediction Scoring System (EPSS).'' [Online]. Available: \url{https://www.first.org/epss/}
\bibitem{cwe} MITRE, ``Common Weakness Enumeration (CWE).'' [Online]. Available: \url{https://cwe.mitre.org/}
\bibitem{cis} Center for Internet Security, ``CIS Amazon Web Services Foundations Benchmark v3.0.0,'' 2024.

\end{thebibliography}

\newcommand{\appchap}[2]{\clearpage\phantomsection\addcontentsline{toc}{chapter}{ภาคผนวก #1 #2}%
  \begin{center}\bfseries\Large ภาคผนวก #1\\[2pt] #2\end{center}\vspace{6pt}}
\appchap{ก}{ตาราง Rubric ของ misconfiguration}\draftnote{}
\appchap{ข}{สคีมา edge list}\draftnote{}
\appchap{ค}{วิธีทำซ้ำผลการทดลอง}\draftnote{}
\appchap{ง}{ความแตกต่างจากข้อเสนอโครงงาน (Deviation log)}\draftnote{GNS3, Nessus, AWS จริง, จำนวนกราฟ}


\end{document}
